\documentclass[11pt]{article}
\usepackage[margin=1in]{geometry}
\usepackage{amsmath,amssymb}
\usepackage{booktabs}
\usepackage{multirow}
\usepackage{xcolor}
\usepackage{caption}
\captionsetup{font=small,labelfont=bf}
\usepackage[colorlinks=true,linkcolor=black,citecolor=black,urlcolor=blue!50!black]{hyperref}
\usepackage{microtype}

\definecolor{court}{HTML}{0E5A3A}
\definecolor{amberx}{HTML}{B26A1F}
\definecolor{steel}{HTML}{5E6E75}
\definecolor{fade}{HTML}{AEB3B5}
\newcommand{\gravity}{\textsc{gravity}}
\newcommand{\ctg}{\textsc{ctg}}
\newcommand{\plainbox}[1]{\par\medskip\noindent\begin{center}%
\colorbox{court!8}{\parbox{0.93\textwidth}{\small \textbf{\textcolor{court}{In plain terms:}} #1}}%
\end{center}\medskip}

% --- publication status: this is self-published, unreviewed research. The footer
% repeats on every page so a saved screenshot of any interior page carries it too.
\usepackage{fancyhdr}
\newcommand{\statusfoot}{%
  \fancyhf{}%
  \fancyfoot[L]{\footnotesize\textcolor{steel}{Working paper · not peer reviewed}}%
  \fancyfoot[C]{\thepage}%
  \fancyfoot[R]{\footnotesize\textcolor{steel}{thegravityreport.com}}%
  \renewcommand{\headrulewidth}{0pt}%
  \renewcommand{\footrulewidth}{0pt}%
}
\fancypagestyle{gravity}{\statusfoot}
\fancypagestyle{plain}{\statusfoot}
\pagestyle{gravity}

\newcommand{\statusnote}{\par\noindent\begin{center}%
\colorbox{fade!25}{\parbox{0.93\textwidth}{\centering\small
\textbf{Working paper --- independent research, not peer reviewed.}\\[2pt]
Findings are provisional and subject to revision. Errors are the author's own.}}%
\end{center}\smallskip}

\title{\textbf{The GRAVITY Specification}\\[4pt]
\large Game-Rate Adjusted Value, Impact, Talent, and Yield\\[2pt]
\normalsize Technical Document · Version 3 · The Gravity Report}
\author{}
\date{July 10, 2026}

\begin{document}
\maketitle
\vspace{-2.5em}
\statusnote

\begin{abstract}
\noindent
This document is the complete, citable specification of \gravity{}, the player-value
system behind ``The 644'' rankings and every numbered research note published at
\texttt{thegravityreport.com}. It defines all inputs, transformations, weights,
adjustments, and boundary rules; the retrospective variant used for historical series;
the rookie module; and the version history. Everything here is sufficient to reproduce
the published numbers given the underlying data. Nothing here is a black box: where the
system makes a judgment call (an exponent, a shrinkage constant, an injury discount),
the call is stated, and the annual Forecast Audit grades what those calls produce.
\end{abstract}

\section{What GRAVITY measures}

\gravity{} asks two questions about every player. \emph{How good is he per minute on the
floor?} --- estimated from box-score impact, scoring efficiency relative to offensive
burden, creation, defensive activity, rebounding, rim pressure, and garbage-time-filtered
lineup impact. \emph{And how much of that ability does a team actually receive?} ---
estimated from games played, sustained role size, and documented injury status. The two
multiply. \gravity{} is therefore a rating of \emph{delivered value}, not pure talent: a
per-minute star who cannot stay on the floor ranks below his talent, on purpose, by a
stated amount.

\plainbox{\gravity{} = (how good per minute) $\times$ (how reliably he's on the floor),
graded against the whole league by the same formula. The scale: 50 = an average NBA
minute, 60+ = a quality starter, 80+ = a franchise engine. The 2025--26 maximum is
Nikola Joki\'c at 97.8.}

\section{Notation and standardization}

Let season $y$ have qualifying player set $\mathcal{P}_y$, with player $i$ playing $m_i$
minutes. For any input statistic $x_k(i)$, define minutes-weighted moments and a clipped
z-score:
\begin{equation}
\mu_k = \frac{\sum_{i \in \mathcal{P}_y} m_i\, x_k(i)}{\sum_i m_i},
\qquad
\sigma_k = \Bigl(\tfrac{\sum_i m_i (x_k(i) - \mu_k)^2}{\sum_i m_i}\Bigr)^{1/2},
\qquad
z_k(i) = \mathrm{clip}\!\left(\frac{x_k(i) - \mu_k}{\sigma_k},\ \pm 4\right).
\label{eq:z}
\end{equation}
Weighting by minutes means $z = 0$ is the \emph{average NBA minute}, not the average
roster spot; replacement level sits near $-1.9$. All standardizations are within-season,
so a $+1\sigma$ season in 2020 and in 2026 mean the same thing relative to their leagues.

\section{Component constructions}

Seven constructed inputs feed the skill blend, each standardized via \eqref{eq:z}:
\begin{align}
\text{eff}(i) &= \bigl(\mathrm{TS}(i) - \mu_{\mathrm{TS}}\bigr) \times 100 \times
\bigl(\mathrm{clip}(\mathrm{USG}(i), 8, 36)/20\bigr)^{0.7}
&&\text{efficiency under load},
\label{eq:eff}\\
\text{creation}(i) &= \mathrm{AST\%}(i) - 0.55\,\mathrm{TOV\%}(i)
&&\text{playmaking net of waste},\\
\text{load}(i) &= \mathrm{USG}(i) + 0.4\,\mathrm{AST\%}(i)
&&\text{shot + playmaking burden},\\
\text{stocks}(i) &= 2.0\,\mathrm{STL\%}(i) + 1.2\,\mathrm{BLK\%}(i)
&&\text{defensive events},\\
\text{rim}(i) &= 100\cdot\mathrm{FTr}(i) + 200\cdot\text{dunk share}(i)
&&\text{rim pressure},\\
\text{vol}(i) &= \text{points per 100 possessions}
&&\text{scoring volume},\\
\text{reb}(i) &= \mathrm{TRB\%}(i)
&&\text{rebounding}.
\end{align}
The exponent $0.7 > 0.5$ in \eqref{eq:eff} is a deliberate design choice: efficiency on a
34\% usage burden is worth more than square-root usage curves grant, and empty-calorie
efficiency at 14\% usage correspondingly less.

\section{Season impact}

The skill blend and the season impact for the current season are
\begin{equation}
S(i) = 0.27\, z_{\text{vol}} + 0.19\, z_{\text{eff}} + 0.22\, z_{\text{crea}}
+ 0.09\, z_{\text{load}} + 0.10\, z_{\text{stk}} + 0.05\, z_{\text{reb}} + 0.08\, z_{\text{rim}},
\label{eq:skill}
\end{equation}
\begin{equation}
I_y(i) = \Bigl[0.44\, z_{\mathrm{BPM}} + 0.12\, z_{\mathrm{WS/48}}
+ 0.14\, z\bigl(\Delta_{\text{on/off}}\cdot \min(1, m_i/2400)\bigr)
+ 0.30\, S(i)\Bigr]\cdot \sqrt{\min(1, m_i/800)},
\label{eq:impact}
\end{equation}
where $\Delta_{\text{on/off}}$ is the garbage-time-filtered on/off differential
(Cleaning the Glass), regressed by minutes before standardization. The trailing factor is
a \emph{within-season} shrink: a scorching 300-minute stretch regresses toward league
average before it can move a ranking. Research Note №2 measured the year-to-year
reliability of filtered on/off at $r = 0.317$ ($n = 184$ players with 1{,}000+ minutes in
consecutive seasons) --- which is why the term carries only 14\% weight and a minutes
regression.

\section{Two-year blend and credibility}

Two seasons blend with minute-proportional weights favoring the recent year, then shrink
toward replacement $\rho = -1.9$ with credibility from actual two-year minutes:
\begin{equation}
w = \frac{0.75\, m_y}{0.75\, m_y + 0.25\, m_{y-1}},
\qquad
R(i) = w\, I_y(i) + (1 - w)\, I_{y-1}(i),
\label{eq:blend}
\end{equation}
\begin{equation}
\tilde{R}(i) = c\,R(i) + (1 - c)\,\rho,
\qquad
c = \sqrt{\min(m_y + m_{y-1},\ 2200)/2200}.
\label{eq:cred}
\end{equation}
An established star with a short injury season keeps his track record; a 300-minute
two-way contract is pulled firmly toward replacement.

\section{Adjustments}

\subsection{Playoff evidence and age}
\begin{equation}
\pi(i) = \mathrm{clip}\Bigl(0.12\,(\mathrm{BPM}^{po} - \mathrm{BPM}^{rs})\cdot
\min(1, m^{po}/300)/2.6,\ \pm 0.30\Bigr),
\label{eq:po}
\end{equation}
\begin{equation}
\alpha(i) = \mathrm{clip}\bigl(0.055\,(24 - \text{age})\,[\text{age} \le 23]
- 0.045\,(\text{age} - 30)\,[\text{age} \ge 31],\ -0.30,\ +0.25\bigr).
\label{eq:age}
\end{equation}

\subsection{Reliability multipliers}
\begin{align}
A(i) &= 0.78 + 0.22\, \min\!\bigl(1,\ (0.7 g_y + 0.3 g_{y-1})/55\bigr)
&&\text{availability},
\label{eq:avail}\\
\Phi(i) &= 0.62 + 0.38\, \min\!\bigl(1,\ \max(\mathrm{mpg}_y,\ 0.85\,\mathrm{mpg}_{y-1})/32\bigr)
&&\text{role size},
\label{eq:role}\\
J(i) &\in [0.78,\ 1.00]
&&\text{documented injury override}.
\label{eq:inj}
\end{align}
Injury overrides are manual, sparse, and published. Table~\ref{tab:inj} lists every
override active in v3.

\begin{table}[h]
\centering\small
\caption{All injury overrides in v3 (July 2026). Nothing else is hand-adjusted.}
\begin{tabular}{llc}
\toprule
Player & Documented reason & $J$ \\
\midrule
Damian Lillard & Achilles tear Apr 2025; missed all 2025--26; age 36 & 0.78 \\
Donte DiVincenzo & Achilles tear 2026 playoffs; out most of 2026--27 & 0.80 \\
Jimmy Butler & ACL tear Jan 19, 2026; out until $\sim$Feb 2027 & 0.80 \\
Tyrese Haliburton & Achilles tear Jun 2025; missed all 2025--26 & 0.82 \\
Kyrie Irving & ACL tear Mar 2025; missed all 2025--26 & 0.88 \\
Fred VanVleet & ACL tear 2025 offseason; missed all 2025--26 & 0.88 \\
Bradley Beal & hip surgery ended 2025--26 after 6 games & 0.90 \\
Dejounte Murray & Achilles rehab; 14 games in 2025--26 & 0.90 \\
Jayson Tatum & Achilles return Mar 2026; strong 16-game sample & 0.95 \\
Ja Morant & injury-marred 20-game 2025--26 (elbow UCL) & 0.95 \\
Joel Embiid & chronic knee management; 38 games & 0.97 \\
\bottomrule
\end{tabular}
\label{tab:inj}
\end{table}

\section{The final scale and tiers}

\begin{equation}
G(i) = 50 + 15\Bigl[\rho + \bigl(\tilde{R} + \pi + \alpha - \rho\bigr)\cdot A \cdot \Phi \cdot J\Bigr].
\label{eq:final}
\end{equation}

\begin{table}[h]
\centering\small
\caption{Tier definitions (fixed score cuts, not fixed counts).}
\begin{tabular}{clc}
\toprule
Tier & Name & \gravity{} score \\
\midrule
T1 & Franchise engines & $\ge 80$ \\
T2 & All-NBA core & 66--80 \\
T3 & All-Star tier & 60--66 \\
T4 & Plus starters & 56--60 \\
T5 & Solid starters & 53--56 \\
T6 & High-leverage rotation & 50--53 \\
T7 & Rotation & 46--50 \\
T8 & Fringe \& developmental & 40--46 \\
T9 & Deep bench, two-ways \& stashes & $< 40$ \\
\bottomrule
\end{tabular}
\label{tab:tiers}
\end{table}

\section{The retrospective variant (GRAVITY-RS)}

Filtered on/off is not uniformly available for early seasons, so historical series
(career arcs, comparables engines) use \gravity-RS: delete the on/off term from
\eqref{eq:impact} and renormalize,
\begin{equation}
I^{RS}_y(i) = \bigl[0.533\, z_{\mathrm{BPM}} + 0.133\, z_{\mathrm{WS/48}}
+ 0.333\, S'(i)\bigr]\cdot \sqrt{\min(1, m_i/800)},
\label{eq:rs}
\end{equation}
with $S'$ using prior-season skill weights
($0.28/0.21/0.22/0.09/0.10/0.06/0.04$; rim term reduced to free-throw rate only).

\section{The rookie module}

Draftees with no NBA minutes enter via a draft-slot prior mapped onto the current
veteran score distribution, adjusted by college production:
\begin{equation}
\text{prod} = 0.5\, z^{c}_{\mathrm{BPM}} + 0.2\, z^{c}_{\mathrm{TS}}
+ 0.15\, z^{c}_{\mathrm{USG}} + 0.15\, z^{c}_{\mathrm{WS/40}},
\qquad
\Delta_{\text{pct}} = \mathrm{clip}(5\,\text{prod},\ \pm 8),
\label{eq:rookie}
\end{equation}
where $z^c$ standardizes within the draft-class cohort. Slot priors (league percentile):
pick 1 $\to$ 84; 2--3 $\to$ 78; 4--7 $\to$ 66; 8--14 $\to$ 56; 15--22 $\to$ 46;
23--30 $\to$ 38; 31--45 $\to$ 28; 46--60 $\to$ 20; minus $0.06$ per pick as a tie-break.

\textbf{Known deficiency, scheduled for v4:} Research Note №1 showed the production blend
\eqref{eq:rookie} would have graded 2019-prospect Ja Morant as ordinary while his college
assist rate sat $+3.5\sigma$ above the 2026 cohort --- creation is the one skill his
collapse never touched. A creation term enters the rookie blend at the next model
revision, before the All-Star 2027 edition.

\section{Version history and edition calendar}

\begin{table}[h]
\centering\small
\caption{Versions are model changes; editions are publications. Old editions are never
revised.}
\begin{tabular}{llp{8.6cm}}
\toprule
& Date & Change \\
\midrule
v1.0 & Jul 10, 2026 & Initial specification; The 644 computed (644 players). \\
v2 & Jul 10, 2026 & Butler injury override added after fact-check (ACL Jan 19, 2026);
his rank moved \#13 $\to$ \#58. Correction published same day. \\
v3 & Jul 10, 2026 & \ctg{} garbage-time-filtered on/off integrated for both blend years
(matched 278 of 279 players with 1{,}000+ minutes); bbref raw on/off retained as
fallback below the match threshold. \\
\midrule
\multicolumn{3}{l}{\textbf{Edition calendar:} Jul 2026 (\textsc{final}) $\to$
All-Star Break Feb 2027 $\to$ End of Season Apr 2027 $\to$} \\
\multicolumn{3}{l}{Post-Playoffs Jun 2027 + Forecast Audit №1 (grades the registered
Queta forecast).} \\
\bottomrule
\end{tabular}
\label{tab:versions}
\end{table}

\section{Data sources and reproducibility}

Season inputs: Basketball-Reference league tables (advanced, per-100-possession,
play-by-play, shooting; regular season and playoffs), 2019--20 through 2025--26. Lineup
inputs: Cleaning the Glass subscriber database (filtered on/off; licensed, not
redistributable --- published outputs are derived values). College inputs:
Sports-Reference CBB. Each edition of The 644 publishes its full output table as CSV
alongside the interactive database. No external player rankings, projections, or draft
boards are consulted at any stage of the pipeline.

\textbf{Known limitations.} Box-score composites undervalue tracking-invisible defense
(the Draymond Green problem, flagged in every affected placement); filtered on/off is
two-thirds noise year-over-year and is weighted accordingly; rookie placements are
priors, not measurements; availability and role discounts are stated choices, and a pure
talent ranking would order some players differently.

\subsection*{Citation}
\begin{quote}
The Gravity Report. \emph{The GRAVITY Specification}, version 3. July 10, 2026.\\
\texttt{https://thegravityreport.com}
\end{quote}

\end{document}
